\chapter{Investigation/Experiment, Result, Analysis and Discussion}
\label{ch:results}

This chapter reports the large-scale evaluation of LectureAssist's question
generator. It describes the environment the experiment was run in, defines every
metric precisely, presents the measured results comparing the agentic pipeline
against the direct-prompting baselines, and discusses what they show --- including a
failure mode of the agentic pipeline that a quality-only evaluation would have
missed entirely.

The experiment is designed around four questions.
\emph{RQ1:} does the agentic pipeline produce better MCQs than single-shot
prompting, including the stronger chain-of-thought and program-of-thought
prompting strategies?
\emph{RQ2:} how does each pipeline behave as the target difficulty rises from
easy to hard?
\emph{RQ3:} is the agent's internal self-validator a reliable quality gate when
checked against an independent judge?
\emph{RQ4:} how trustworthy is the LLM judge itself --- checked by calibrating
the answer-solving model on questions with known answers, so that agreement
between models is never mistaken for correctness?

\section{Environment Setup}
\label{sec:res-env}

All experiments were run on a single Google Colab instance with an NVIDIA T4 GPU
($16$\,GB), with models served locally by Ollama. No paid or proprietary API is used
anywhere in the \emph{generation} path --- every question in the study is produced
entirely locally. (Cross-family judging additionally re-runs the \emph{measurement}
through a cloud API; Section~\ref{sec:res-metrics} explains why this does not change
the system under evaluation.) Bulk-evaluation artifacts are written to Google Drive
so that a run survives a Colab runtime disconnection and can be resumed rather than
restarted.

The generator is \texttt{gemma3:4b}, pinned: it is never allowed to escalate to the
larger model, even when it returns malformed JSON, so that ``generated by the $4$B''
is true of every question in the study. The judge and the Answer Generator are
\texttt{gemma3:12b}. The model that scores or solves a question is therefore never
the model that wrote it.

\paragraph{Source data.} Questions are generated from a college-level STEM document
with a clean text layer --- the \emph{Limits} chapter of OpenStax \emph{Calculus,
Volume~1}~\cite{openstax_calculus} --- ingested through the standard document
pipeline (extraction $\rightarrow$ map-reduce summary $\rightarrow$ retrieval
index). The extracted chapter is $92{,}447$ characters over $71$ pages. Every
pipeline generates from the identical extracted content and reads the same number of
content windows, so any difference in the output reflects the generation method and
not the input (Section~\ref{sec:meth-controls}).

\paragraph{Calibration set.} A separate, stratified set of $148$ MCQ with known
answers is used to calibrate the Answer Generator: $100$ \emph{computational} items
whose answers are computed with SymPy, and $48$ \emph{conceptual} items drawn from
the chapter's definitions and theorems. This set never enters the generation path.

\section{Testing and Evaluation}
\label{sec:res-metrics}

Each question set is measured along four metric families. Let $Q = \{q_1,\dots,q_n\}$
be a generated set and $C = \{c_1,\dots,c_m\}$ the lecture chunks.

\paragraph{Evaluation harness.} Generation and measurement are deliberately
decoupled. The bulk run in Colab only \emph{generates}: all questions, together
with the self-validator's per-question verdicts, are saved to a single JSON
artifact. Judging is then performed by a standalone harness that reads this
artifact and scores every question individually. The harness exists as two twin
scripts with identical rubric, solver, metrics and outputs, differing
\emph{only} in the model backend: one judges with the local Ollama model
(\texttt{gemma3:12b}); the other with any OpenAI-compatible API, defaulting to
DeepSeek-R1 --- the judge model EQGBench itself uses --- so the same corpus can
be measured by a judge from a different model family. The two backends are kept
as separate, side-by-side programs precisely so that the local run and the
cross-family run remain directly comparable and neither can silently change the
other; pointing both at the same questions yields inter-judge agreement between
a Gemma judge and a DeepSeek judge. Judged rows are appended to a resumable
per-question log, so an interrupted judging run continues instead of
restarting, and every judge call fails closed. The cloud backend is a
\emph{measurement instrument} only: the generator under evaluation is still the
local \texttt{gemma3:4b}, and measuring a local system with a cloud instrument
does not make the deployed system non-local.

\paragraph{Relation to EQGBench.} The protocol follows the evaluation frame of
EQGBench~\cite{eqgbench}, a benchmark for LLM educational question generation
that scores output with an LLM judge along per-question quality dimensions plus
an acceptance decision, rather than with $n$-gram overlap. Overlap metrics
(BLEU/ROUGE) are deliberately not reported: generation here is open-ended over a
whole chapter, so there is no reference question set to overlap with, and
similarity to any particular phrasing is not evidence of quality. EQGBench's
own results also motivate a design choice in this chapter: in that benchmark
most quality dimensions saturate near the ceiling for every model while a
single demanding dimension (competence-oriented guidance) collapses and does
the discriminating. The same pattern appears here --- judge scores compress
near the top of the scale (Table~\ref{tab:bulk-quality}) --- which is why this
evaluation adds repetition- and grounding-based measures (duplicate rate,
diversity, effective accept) that EQGBench omits, and lets them carry the
discriminating signal.

\paragraph{Judge metrics.} An independent judge model scores each question from $1$
to $5$ on \emph{correctness} (is it factually right and is the marked option truly
the best one), \emph{clarity} (is the stem unambiguous), \emph{difficulty match}
(does it match the difficulty it was asked for), \emph{distractor quality} (are the
wrong options plausible but decidably wrong), and an \emph{overall} score; it also
returns a binary \texttt{accept}. The \textbf{accept rate} is
\begin{equation}
\mathrm{AcceptRate}(Q) \;=\; \frac{1}{n}\sum_{i=1}^{n}\mathbb{1}[\texttt{accept}(q_i)] .
\end{equation}
A judging call that errors or returns unparseable output \emph{fails closed} --- it
is scored $0$ and rejected, never quietly dropped, so a fragile pipeline cannot
improve its own score by crashing the judge. An \texttt{overall} score below $3$
forces \texttt{accept} to false for consistency. The exact rubric sent to the
judge, shown below, instructs it to use the full $1$--$5$ range --- an
anti-saturation measure whose necessity Section~\ref{sec:res-findings}
confirms:

\begin{verbatim}
You are an impartial, strict but fair exam-quality judge. Score this
ONE MCQ question on a 1-5 scale for each dimension.

Target difficulty: [DIFF] - [difficulty description]
Content snippet (ground truth):
[lecture content]

Question:
Q: [question text]      Options: [options]      A: [answer]

Score 1 (worst) to 5 (best) on:
1. correctness        - factually correct, answerable from content
2. clarity            - unambiguous, well-formed wording
3. difficulty_match   - actual difficulty matches the target
4. distractor_quality - wrong options plausible but clearly wrong
5. overall            - holistic quality

BE DISCRIMINATING - use the FULL range, do not default to 5:
  5 = flawless; 4 = good; 3 = ordinary; 2 = weak; 1 = poor.
Then decide accept (true) or reject (false).

Return STRICT JSON:
{"correctness":1,"clarity":1,"difficulty_match":1,
 "distractor_quality":1,"overall":1,"accept":true,"reason":""}
\end{verbatim}

\paragraph{Deterministic metrics.} Two questions are duplicates if their stems have
a character-level similarity ratio above $0.85$. The \textbf{duplicate rate} is the
fraction of questions that duplicate an earlier question in the set. Since a
pipeline can trivially achieve a perfect accept rate by asking the same good
question fifty times, quality and repetition are combined into the
\textbf{effective accept rate}, which is the metric used for ranking:
\begin{equation}
\label{eq:effacc}
\mathrm{EffAccept}(Q) \;=\; \mathrm{AcceptRate}(Q)\;\times\;\bigl(1 - \mathrm{DupRate}(Q)\bigr) .
\end{equation}

\paragraph{Embedding metrics.} Using the same Sentence-BERT encoder as the retrieval
engine, \textbf{grounding} is the mean best-match cosine similarity of each question
to any lecture chunk, and \textbf{diversity} is one minus the mean
nearest-neighbour cosine similarity between questions:
\begin{equation}
\mathrm{Ground}(Q) = \frac{1}{n}\sum_{i=1}^{n}\max_{j}\ \mathrm{sim}(q_i, c_j),
\qquad
\mathrm{Div}(Q) = 1 - \frac{1}{n}\sum_{i=1}^{n}\max_{j \neq i}\ \mathrm{sim}(q_i, q_j).
\end{equation}
These require no LLM at all, which makes them a useful check on the judge: they
cannot be talked into a high score by fluent writing.

\paragraph{Answer metrics.} The \textbf{independent-match rate} is the fraction of
questions where the Answer Generator, solving the question \emph{without seeing} the
marked option, chooses the same option the generator marked. The
\textbf{judge-verified rate} is the fraction the judge confirms as correctly marked.

\paragraph{Validator--judge agreement.} For the agentic pipeline only, the
self-validator's \texttt{PASS} is treated as a prediction and the independent judge's
\texttt{accept} as ground truth, and the $F_1$ score between them is reported. This
measures whether the agentic pipeline's own quality control is actually detecting
what an outside judge would detect, or is merely rubber-stamping its own output.

\paragraph{Judge reliability.} The generator and the judge share the Gemma-3
family, so they may share blind spots, and an evaluation that trusted the judge
blindly would inherit them. Two safeguards address this. First, the solver
calibration of Table~\ref{tab:bulk-gold} measures the judging-side models on
questions with \emph{known} answers, which is what allows
Section~\ref{sec:res-findings} to bound how much the agreement-based answer
metrics can be trusted rather than taking them at face value. Second, the twin
cross-family judge described under \emph{Evaluation harness} above re-measures
the identical corpus with DeepSeek-R1 and reports inter-judge agreement.
Cross-family re-judging is part of the full four-pipeline run whose CoT and PoT
cells are pending in Tables~\ref{tab:bulk-quality} and~\ref{tab:bulk-answer};
the results reported in this chapter are from the local \texttt{gemma3:12b}
judge alone.

\section{Results and Discussion}
\label{sec:res-results}

Table~\ref{tab:bulk-quality} reports question quality and
Table~\ref{tab:bulk-answer} reports diversity, grounding and answer correctness, in
both cases per difficulty and per pipeline. Table~\ref{tab:bulk-gold} reports the
Answer Generator calibration and Table~\ref{tab:bulk-f1} the validator--judge
agreement.

\begin{table}[H]
\centering
\caption{Question quality by difficulty and pipeline (judge scores $1$--$5$; accept
rate in \%). CoT and PoT rows are pending the full four-pipeline run and are not yet
measured.}
\label{tab:bulk-quality}
\small
\begin{tabular}{|l|l|c|c|c|c|c|}
\hline
\textbf{Difficulty} & \textbf{Pipeline} & \textbf{Overall} & \textbf{Correctness} &
\textbf{Clarity} & \textbf{Diff.\ match} & \textbf{Accept} \\
\hline
Easy   & Agentic     & 4.956 & 4.94 & 4.97 & 4.60 & 100\% \\
       & Non-agentic & 4.854 & 4.88 & 4.90 & 4.40 & 100\% \\
       & CoT         & ---   & ---  & ---  & ---  & ---   \\
       & PoT         & ---   & ---  & ---  & ---  & ---   \\
\hline
Medium & Agentic     & 4.939 & 4.90 & 4.95 & 4.45 & 99\% \\
       & Non-agentic & 4.814 & 4.80 & 4.85 & 4.20 & 97\% \\
       & CoT         & ---   & ---  & ---  & ---  & ---  \\
       & PoT         & ---   & ---  & ---  & ---  & ---  \\
\hline
Hard   & Agentic     & 4.884 & 4.88 & 4.90 & 4.20 & 97\% \\
       & Non-agentic & 4.788 & 4.60 & 4.80 & 4.00 & 91\% \\
       & CoT         & ---   & ---  & ---  & ---  & ---  \\
       & PoT         & ---   & ---  & ---  & ---  & ---  \\
\hline
\end{tabular}
\end{table}

\begin{table}[H]
\centering
\caption{Diversity, grounding and answer correctness by difficulty and pipeline.
CoT and PoT rows are pending the full four-pipeline run and are not yet measured.}
\label{tab:bulk-answer}
\footnotesize
\setlength{\tabcolsep}{4pt}
\begin{tabular}{|l|l|c|c|c|c|c|c|}
\hline
\textbf{Difficulty} & \textbf{Pipeline} & \textbf{Duplicate} & \textbf{Grounding} &
\textbf{Diversity} & \textbf{Eff.\ accept} & \textbf{Indep.\ match} &
\textbf{Judge-verified} \\
\hline
Easy   & Agentic     & 35\% & 0.562 & 0.74 & 65\% & 97\% & 100\% \\
       & Non-agentic & 63\% & 0.642 & 0.55 & 37\% & 95\% & 99\%  \\
       & CoT         & ---  & ---   & ---  & ---  & ---  & ---   \\
       & PoT         & ---  & ---   & ---  & ---  & ---  & ---   \\
\hline
Medium & Agentic     & 38\% & 0.595 & 0.70 & 62\% & 88\% & 99\% \\
       & Non-agentic & 70\% & 0.614 & 0.48 & 30\% & 84\% & 95\% \\
       & CoT         & ---  & ---   & ---  & ---  & ---  & ---  \\
       & PoT         & ---  & ---   & ---  & ---  & ---  & ---  \\
\hline
Hard   & Agentic     & 70\% & 0.513 & 0.45 & 30\% & 88\% & 99\% \\
       & Non-agentic & 38\% & 0.616 & 0.60 & 62\% & 83\% & 81\% \\
       & CoT         & ---  & ---   & ---  & ---  & ---  & ---  \\
       & PoT         & ---  & ---   & ---  & ---  & ---  & ---  \\
\hline
\end{tabular}
\end{table}

\begin{table}[H]
\centering
\caption{Answer-Generator calibration: accuracy on the known-answer set. This is a
calibration of the \emph{solver}, not a score on the generated questions.}
\label{tab:bulk-gold}
\small
\begin{tabular}{|l|l|c|}
\hline
\textbf{Model} & \textbf{Role} & \textbf{Accuracy on known-answer set} \\
\hline
\texttt{gemma3:4b}  & Generator model      & 70\% \\
\texttt{gemma3:12b} & Judge / solver model & 55\% \\
\hline
\end{tabular}
\end{table}

\begin{table}[H]
\centering
\caption{Validator--judge agreement for the agentic pipeline: does the pipeline's own
self-validator detect what an independent judge detects?}
\label{tab:bulk-f1}
\small
\begin{tabular}{|l|c|}
\hline
\textbf{Difficulty} & $\mathbf{F_1}$ \\
\hline
Easy   & 0.786 \\
Medium & 0.774 \\
Hard   & 0.842 \\
\hline
\end{tabular}
\end{table}

\subsection{Findings}
\label{sec:res-findings}

\paragraph{1. The agentic pipeline produces more answer-correct questions, and the
gap widens with difficulty.} Judge-verified correctness is $99\%$ for the agentic
pipeline versus $81\%$ for the plain baseline on hard questions, and the independent
re-solve rate is higher at every difficulty ($97/88/88\%$ against $95/84/83\%$). Mean
judge score and accept rate favour the agentic pipeline throughout. This is the
result the architecture was built for: the validation-and-retry loop catches
questions whose marked answer is wrong, and it matters most exactly where the $4$B
generator is weakest. Note that the easy and medium gaps are small --- the plain
baseline is already near-perfect there --- so the agentic machinery buys little on
easy material and a great deal on hard material.

\paragraph{2. Diversity, not quality, is the discriminator --- and it reverses at
hard.} Judge scores are compressed into a narrow band near the ceiling ($4.79$ to
$4.96$ out of $5$), which is a known property of LLM judges and is why this
evaluation does not rely on them alone. The duplicate rate separates the pipelines
far more sharply than any judge dimension. At easy and medium the agentic pipeline is
much less repetitive ($35\%$ and $38\%$ duplicates against $63\%$ and $70\%$), but at
hard it \emph{regresses badly} to $70\%$ duplicates against the baseline's $38\%$.
Because the effective accept rate (Eq.~\ref{eq:effacc}) discounts repetition, this
single metric flips the ranking at the hard level: $30\%$ for the agentic pipeline
against $62\%$ for the baseline.

The mechanism is visible in the grounding column. Agentic grounding \emph{falls} at
hard ($0.513$, its lowest value) while the baseline's stays flat ($0.616$). Asked for
hard questions, the planner and validator converge on the narrow band of the chapter
that actually supports difficult questions, retrieve repeatedly from it, and then ask
about it over and over. The retry loop makes this worse rather than better: each
retry is another draw from the same small region. This is a genuine defect of the
current design, it is reported rather than hidden, and it is the clearest direction
for future work (Section~\ref{sec:concl-future}).

\paragraph{3. The Answer Generator is a weak instrument, and this bounds the answer
metrics.} On the known-answer set, the independent solver (\texttt{gemma3:12b})
scores $55\%$ and the generator model (\texttt{gemma3:4b}) scores $70\%$ --- neither
is close to reliable, and the larger model is, counter-intuitively, the worse solver
on this material. This matters for interpretation. The high generator--solver
agreement ($88$--$97\%$) therefore cannot be read as $88$--$97\%$ of questions being
verifiably correct: two models from the same family, agreeing at a rate far above the
accuracy either achieves on questions with known answers, are partly agreeing because
they share the same behaviour, not because they are both right. The
independent-match rate is a \emph{relative} signal between pipelines and must be read
alongside this calibration. Reporting it without the calibration would have made the
system look far more trustworthy than the evidence supports.

\paragraph{4. Self-validation works, but imperfectly.} The agentic pipeline's own
validator agrees with the independent judge at $F_1 = 0.77$ to $0.84$. It is
therefore doing real work --- it is not rubber-stamping its own output --- but roughly
one in five of its verdicts would not survive external review, which is a useful
bound on how much the pipeline's self-reported accept rate can be trusted.

\paragraph{5. The agentic pipeline is roughly seven times more expensive.} Generating
$50$ easy questions took about $28$ minutes agentically against about $4$ minutes for
the plain baseline, since each question costs a plan lookup, one or more retrievals,
a generation call, a validation call, and possibly several retries. That cost is the
price of the correctness and diversity gains at easy and medium difficulty --- and at
hard difficulty, where the effective accept rate reverses, it currently buys nothing.

\section{Summary}
\label{sec:res-summary}

The evaluation supports a qualified conclusion. The agentic pipeline generates
questions that are better judged and, more importantly, more likely to be marked with
the \emph{correct} answer --- decisively so on hard questions, which is where an
automatic quiz generator is most likely to mislead a student. It does so at roughly
seven times the wall-clock cost, and it has a real weakness: on hard material it
collapses onto a narrow slice of the chapter and repeats itself, to the point that a
repetition-aware metric ranks the plain baseline above it. Two methodological points
carry beyond this system. First, a question generator evaluated on quality alone will
look excellent regardless of whether it marks the right answer, so answer correctness
must be measured separately, by a model that does not see the marked option. Second,
an LLM judge must itself be calibrated against questions with known answers, or its
verdicts will be over-trusted; here that calibration is what prevented a $97\%$
agreement rate from being reported as $97\%$ correctness.
